\section{A taxonomy of the field, and when two studies are comparable}\label{sec:taxonomy}

The literature labels a study ``generalized thermoelastic'' or it does not. The single label collapses several independent commitments into one word. A study may adopt a relaxation law for the heat flux, a gradient or
integral law for the stress, a coupling to a second field, and a validation regime. Each of these can
be varied while the others are held fixed. Two studies that share the label can therefore differ in the
mathematics of their governing equations, in their modal content and in the evidence they offer.
\Cref{sec:synthesis} compares families; this section says what a family \emph{is} as a coordinate, so that a
new model can be placed without argument.

The heat-flux template of \cref{sec:foundations} is a special case of Axis~2. Its flux law is local in
space and first order in the flux, with an optional thermal-displacement ($1/s$) term, and the lags are the coordinate that separates the template's members. A
model that needs an integral, fractional or gradient operator has left that axis rather than moved along
it. Such models therefore cannot be ranked against the template's members by setting a lag to zero.

The three groups of \cref{tab:taxonomy} (structure of the heat law, elastic constitutive law, additional fields) are a coarse, model-oriented view. In the five-axis scheme below, these commitments correspond mainly to the mathematical axis (Axis~2), both for the heat law and for the local, nonlocal or gradient form of the elastic law. The additional fields correspond to the physical and coupling axes (Axes~1 and 4). The wave and validation axes (Axes~3 and 5) describe what a study does with the model.

A study is classified by its vector of commitments $\mathbf{c}=(c_1,\dots,c_5)$, one coordinate per axis:

\begin{enumerate}[label=\textbf{Axis \arabic*:}]
\item \textbf{Physical axis (material class).} Classical elastic, thermoelastic, piezoelectric,
  piezothermoelastic, magneto-thermo-elastic, poroelastic, micropolar, two-temperature, semiconductor-like.
  The physical axis fixes which conserved quantities and driving fields exist, before any approximation is made.
\item \textbf{Mathematical axis (operator structure).} Local PDE, integro-differential, fractional-order,
  nonlocal gradient, hereditary memory integral, or variational/weak form. The most consequential
  differences between studies run along this axis.
\item \textbf{Wave axis (modal content).} Does the thermal mode propagate at finite speed; how many coupled
  modes exist; is there attenuation; what is the dispersion character of each.
\item \textbf{Coupling axis (field topology).} Which fields are coupled and how strongly:
  thermoelastic, thermo-piezoelectric, thermo-magneto-elastic, thermo-poro-elastic, thermo-micro-elastic.
  Coupling topology, and not the choice of energy functional, decides whether an interface problem is
  tractable.
\item \textbf{Validation axis (evidence regime).} Existence and uniqueness only; closed-form solution for a
  special geometry; numerical solution validated against a reference; comparison with experiment; a shared
  benchmark suite.
\end{enumerate}

Two studies are \emph{structurally} comparable, at the level of model class, when their vectors agree in the first two coordinates. A comparison of numbers further needs compatible geometry, parameters, boundary and interface conditions, nondimensional regime and observable. If they
differ there, a comparison of numbers means nothing and only a comparison of structure is available. To
make the weaker statement quantitative, define the axis-weight dissimilarity
\begin{equation}\label{eq:taxdissim}
d(\mathbf{c},\mathbf{c}')=\sum_{j=1}^{5} w_j\,\bigl(1-\delta_{c_j c_j'}\bigr),
\qquad
w_{\mathrm{math}}\ge w_{\mathrm{phys}}>w_{\mathrm{coup}}>w_{\mathrm{wave}}>w_{\mathrm{val}} ,
\end{equation}
with the ordering stated because it is a modeling decision open to challenge. A difference in the mathematical axis changes the governing operator and invalidates most direct comparisons. A difference in the validation axis changes only the epistemic status of the conclusion. The weights are reported whenever the dissimilarity is used, and no result in this review depends on their numerical values.

\begin{table}[H]\centering\small
\caption{The five axes with their admissible values. A study occupies one cell of the product; the last
column gives an example assignment for a coupled fractional piezothermoelastic problem.}
\label{tab:taxonomy-axes}
\begin{tabularx}{\textwidth}{@{}p{2.1cm}Y Y@{}}
\toprule
\textbf{Axis} & \textbf{Admissible values} & \textbf{Example} \\
\midrule
Physical & elastic; thermoelastic; piezoelectric; piezothermoelastic; magneto-thermo-elastic; poroelastic;
  micropolar; two-temperature; semiconductor & piezothermoelastic \\
Mathematical & local PDE; integro-differential; fractional; nonlocal gradient; hereditary; variational &
  fractional + local coupling \\
Wave & finite/infinite thermal speed; number of coupled modes; attenuating/non-attenuating &
  two modes, one attenuated \\
Coupling & none; thermal; thermal + electric; thermal + magnetic; thermal + poro; thermal + micro-rotation &
  thermal + electric \\
Validation & proof only; closed form; verified numerics; experiment; shared benchmark &
  verified numerics \\
\bottomrule
\end{tabularx}
\end{table}

A taxonomy earns its place only if it changes a decision, and this one changes three decisions in this review.
The taxonomy fixes the order in which \cref{sec:synthesis} compares families. It is also why \cref{sec:interfaces} insists that an energy flux be defined consistently with the model before ratios are compared (see \cref{sec:eqregister}). It is also why \cref{sec:open} keeps apart the questions that differ in the
mathematical axis from those that differ only in evidence.
